\documentclass[a4paper,12pt]{article}
\usepackage{color}
\usepackage{graphicx, gensymb}
\usepackage{amsfonts, cancel, amssymb}
\usepackage{amsmath, amsthm, array, esvect, esint}
\usepackage{typearea}
\usepackage[a4paper,left=2cm,right=2cm,top=2cm,bottom=3cm,bindingoffset=5mm]{geometry}
\newcommand\numberthis{\addtocounter{equation}{1}\tag{\theequation}}
\newcommand{\bra}{}
\newcommand{\kom}{}
\newcommand{\brak}{}
\newcommand{\braket}{}
\newcommand{\intx}{}
\newcommand{\intr}{}
\newcommand{\kla}{}
\newcommand{\klam}{}
\newcommand{\klammer}{}
\newcommand{\oper}{}
\newcommand{\tecxt}{}
\newcommand{\Bra}{}
\newcommand{\Ket}{}

\begin{document}

\title{Einstein field equations}
\author{Joachim Schwardt}
\date{\today}
\maketitle

\section{Newtonian gravity}
According to \textsc{Newton}ian gravity, there is an attractive force between any two massive objects. Let $G$ be the gravitational constant and $\vec{r}_m$ the position of $m$, while $M$ sits at the center of the coordinate frame:
\begin{align*}
\vec{F}_m&=-GmM\frac{\vec{r}_m}{r_m^3}=m\ddot{\vec{r}}_m\numberthis\label{Newton force}
\end{align*}
We can rewrite (\ref{Newton force}) as a field equation by factoring out $m$ on both sides and defining the gravitational field generated by $M$ as $\vec{g}:=\ddot{\vec{r}}_m$:
\begin{align*}
\vec{g}&=-GM\frac{\vec{r}}{r^3}\numberthis\label{Newton pointlike field equation}
\end{align*}
The total mass $M$ can be expressed as an integral over the matter density $\rho(\vec{r})$. Also note that the right-hand side of (\ref{Newton pointlike field equation}) can be written as a gradient of a scalar field, which means that we can define $\vec{g}:=\nabla\Phi$:
\begin{align*}
\vec{g}=-G\intr\int_{\mathbb{R}^{3}}{\rho(\vec{s})\frac{\vec{r}-\vec{s}}{|\vec{r}-\vec{s}|^3}}\,\mathrm{d}^{3}s=-G\nabla\intr\int_{\mathbb{R}^{3}}{\frac{\rho(\vec{s})}{|\vec{r}-\vec{s}|}}\,\mathrm{d}^{3}s=\nabla\Phi \numberthis\label{Newton integral representation}
\end{align*}
Note that $\nabla^2\frac{1}{r}=-4\pi\delta(r)$, which gives us the \textsc{Newton}ian gravitational field equation:
\begin{align*}
\nabla^2\Phi=-G\intr\int_{\mathbb{R}^{3}}{\rho(\vec{s})\nabla^2\frac{1}{|\vec{r}-\vec{s}|}}\,\mathrm{d}^{3}s=4\pi G\rho\numberthis\label{Newton field equation}
\end{align*}
\section{Introduction to tensor calculus}
\subsection{Definion and transformation of tensors}[\footnote{Uni Tuebingen Kokkotas GTR files}]

%https://www.tat.physik.uni-tuebingen.de/~kokkotas/Teaching/GTR_files/GTR2019_1.pdf

Let $x^\mu$, $x^{\tilde{\mu}}$ be two coordinate systems, where these are connected by equations of the form $x^{\tilde{\mu}}=x^{\tilde{\mu}}(x^\nu)$ and $x^\mu=x^\mu(x^{\tilde{\nu}})$. A manifold is called differential, if $A^{\tilde{\mu}}_\nu:=\frac{\partial x^{\tilde{\mu}} }{\partial x^\nu}$ and $A^\mu_{\tilde{\nu}}:=\frac{\partial x^\mu }{\partial x^{\tilde{\nu}}}$ exist. The transformation of any contravariant vector $v^\mu\rightarrow v^{\tilde{\mu}}$ is then given by:
\begin{align*}
v^\mu&=\sum_{\tilde{\nu}} A^\mu_{\tilde{\nu}} v^{\tilde{\nu}}\equiv A^\mu_{\tilde{\nu}} v^{\tilde{\nu}}\numberthis\label{Contravariant transformation}
\end{align*}
The transformation of a covariant vector $v_\mu\rightarrow\tilde{v}_\mu$ is on the other hand given by:
\begin{align*}
v_\mu&=\sum_{\tilde{\nu}} A^{\tilde{\nu}}_\mu v_{\tilde{\nu}}\equiv A^{\tilde{\nu}}_\mu v_{\tilde{\nu}}\numberthis\label{Covariant transformation}
\end{align*}
A tensor of rank 2 is a quantity $T^{\mu\nu}$, such that for any $a_\mu,b_\nu$ and any coordinates $x^\mu\rightarrow x^{\tilde{\mu}}$, the following scalar is form-invariant:
\begin{align*}
T^{\mu\nu}a_\mu b_\nu&=T^{\tilde{\mu}\tilde{\nu}}a_{\tilde{\mu}}b_{\tilde{\nu}} =:\phi\numberthis\label{Tensor rank 2 definition}
\end{align*}
The transformations of this tensor are then given by:
\begin{align*}
T^{\mu\nu}&=A^\mu_{\tilde{\alpha}} A^\nu_{\tilde{\beta}} T^{\tilde{\alpha}\tilde{\beta}}&T^\mu_\nu&=A^\mu_{\tilde{\alpha}} A^{\tilde{\beta}}_\nu T^{\tilde{\alpha}}_{\tilde{\beta}}&T_{\mu\nu}&=A^{\tilde{\alpha}}_\mu A^{\tilde{\beta}}_\nu T_{\tilde{\alpha}\tilde{\beta}}\numberthis\label{Tensor rank 2 transformation}
\end{align*}
\subsection{Differentials of tensors}
Define the abbreviations $\phi_{,\alpha}\equiv\partial_\alpha\phi:=\frac{\partial\phi}{\partial x_\alpha}$ and $\phi_{,\tilde{\alpha}}\equiv\partial_{\tilde{\alpha}}\phi:=\frac{\partial\phi}{\partial \tilde{x}_\alpha}$. Note that the previously defined transformation matrices are then written as $A^\mu_{\tilde{\nu}}=x^\mu_{,\tilde{\nu}}$ and $A^{\tilde{\mu}}_\nu=x^{\tilde{\mu}}_{,\nu}$. Then the derivative of a scalar field $\phi$ is:
\begin{align*}
\phi_{,\alpha}&=x^{\tilde{\beta}}_{,\alpha}\phi_{,\tilde{\beta}}\numberthis\label{Derivative scalar field}
\end{align*}
Using $v^{\tilde{\nu}}=x^{\tilde{\nu}}_{,\kappa}v^\kappa$ in the last step, the derivative of a contravariant vector $v^\mu$ is given by:
\begin{align*}
v^\mu_{,\alpha}&=(x^\mu_{,\tilde{\nu}}v^{\tilde{\nu}})_{,\alpha}=x^\mu_{,\alpha\tilde{\nu}}v^{\tilde{\nu}}+x^\mu_{,\tilde{\nu}}v^{\tilde{\nu}}_{,\alpha}=x^\mu_{,\tilde{\rho}\tilde{\nu}}x^{\tilde{\rho}}_{,\alpha}x^{\tilde{\nu}}_{,\kappa}v^\kappa+x^\mu_{,\tilde{\nu}}x^{\tilde{\rho}}_{,\alpha}v^{\tilde{\nu}}_{,\tilde{\rho}}\numberthis\label{Derivative contravariant vector step 1}
\end{align*}
When we define the christoffel symbol $\Gamma^\mu_{\alpha\kappa}:=x^\mu_{,\tilde{\rho}\tilde{\nu}}x^{\tilde{\rho}}_{,\alpha}x^{\tilde{\nu}}_{,\kappa}$, we can rewrite (\ref{Derivative contravariant vector step 1}) in such a way, that we get a new quantity, that once again transforms like a tensor of rank 2:
\begin{align*}
v^\mu_{,\alpha}-\Gamma^\mu_{\alpha\kappa}v^\kappa&=x^\mu_{,\tilde{\nu}}x^{\tilde{\rho}}_{,\alpha}v^{\tilde{\nu}}_{,\tilde{\rho}}\numberthis\label{Derivative contravariant vector step 2}
\end{align*}
Finally, the absolute derivative of a contravariant vector is defined as:
\begin{align*}
v^\mu_{;\alpha}&:=v^\mu_{,\alpha}+\Gamma^\mu_{\alpha\kappa}v^\kappa\numberthis\label{Absolute derivative of contravariant vector}
\end{align*}
The absolute derivatives of other tensors are:
\begin{align*}
\phi_{;\alpha}&=\phi_{,\alpha}\\
v_{\mu ;\alpha}&=v_{\mu ,\alpha}-\Gamma^\kappa_{\alpha\mu}v_\kappa\\
T^{\mu\nu}_{;\alpha}&=T^{\mu\nu}_{,\alpha}+\Gamma^\mu_{\kappa\alpha} T^{\kappa\nu}+\Gamma^\nu_{\kappa\alpha}T^{\mu\kappa}\\
T^\mu_{\nu ;\alpha}&=T^\mu_{\nu ,\alpha}+\Gamma^\mu_{\kappa\alpha}T^\kappa_\nu -\Gamma^\kappa_{\nu\alpha}T^\mu_\kappa
\end{align*}
\subsection{Parallel transport and the curvature tensor}
Let $B$ and $C$ be two neighboring points with displacement $dx^\mu$ and $a_\mu(B)$ a covariant vector. Then we get:
\begin{align*}
a_\mu(C)&=a_\mu(B)+a_{\mu,\nu}(B)dx^\nu
\end{align*}
The parallel transport of the vector however, is to be considered equivalent to $a_\mu(B)$ in a certain way. We write $A_\mu(C):=a_\mu(B)+\delta a_\mu(B)$, where:
\begin{align*}
\delta a_\mu &=\Gamma^\kappa_{\mu\nu}a_\kappa dx^\nu\numberthis\label{Parallel transport of covariant vector}\\
\delta a^\mu &=-\Gamma^\mu_{\kappa\nu}a^\kappa dx^\nu\numberthis\label{Parallel transport of contravariant vector}
\end{align*}
Note that $\delta\phi=0$. Consider now the effect of a second parallel transport from $C$ to $D$ (quantities are given at the point $B$, unless stated otherwise):
\begin{align*}
a^\mu(C)&=a^\mu-\Gamma^\mu_{\kappa\nu}a^\kappa dx^\nu\numberthis\label{Parallel transport B to C}\\
a^\mu(D)&=a^\mu(C)-\Gamma^\mu_{\rho\lambda}(C)a^\rho(C)\delta x^\lambda=\\
&=a^\mu-\Gamma^\mu_{\kappa\nu}a^\kappa dx^\nu-\Gamma^\mu_{\rho\lambda}(C)[a^\rho -\Gamma^\rho_{\kappa\nu}a^\kappa dx^\nu ]\delta x^\lambda\numberthis\label{Parallel transport C to D}
\end{align*}
Using the approximation $\Gamma^\mu_{\rho\lambda}(C)\simeq \Gamma^\mu_{\rho\lambda}+\Gamma^\mu_{\rho\lambda,\alpha}dx^\alpha$, we can express (\ref{Parallel transport C to D}) with quantities at $B$ only. In the spirit of a second order approximation (curvature!), we neglect the term with $(dx)^3$:
\begin{align*}
a^\mu(D)&\simeq a^\mu-\Gamma^\mu_{\kappa\nu}a^\kappa dx^\nu-(\Gamma^\mu_{\rho\lambda}+\Gamma^\mu_{\rho\lambda,\alpha}dx^\alpha)[a^\rho -\Gamma^\rho_{\kappa\nu}a^\kappa dx^\nu ]\delta x^\lambda\simeq\\
&\simeq a^\mu-\Gamma^\mu_{\kappa\nu}a^\kappa dx^\nu-\Gamma^\mu_{\rho\lambda}a^\rho\delta x^\lambda-\Gamma^\mu_{\rho\lambda,\alpha}a^\rho dx^\alpha\delta x^\lambda +\Gamma^\mu_{\rho\lambda}\Gamma^\rho_{\kappa\nu}a^\kappa dx^\nu\delta x^\lambda
\end{align*}
If we define a point $C'$, such that the four points form a parallelgram with sides $dx^\mu$ and $\delta x^\mu$, then we can also get $a^\mu(D)'$ via the path $B\rightarrow C'\rightarrow D$. This will only swap $d$ and $\delta$ in the respective cases, leaving $dx^\alpha$ unchanged:
\begin{align*}
\delta a^\mu &:=a^\mu(D)-a^\mu(D)'=(\Gamma^\mu_{\kappa\lambda}\Gamma^\kappa_{\rho\alpha}+\Gamma^\mu_{\rho\lambda,\alpha})a^\rho (dx^\alpha\delta x^\lambda -\delta x^\alpha dx^\lambda)\numberthis\label{Parallel transport B to D}
\end{align*}
Swapping the indices $\alpha$ and $\lambda$ in (\ref{Parallel transport B to D}) and adding the both equations gives the following result, where $R^\mu_{\rho\lambda\alpha}$ is called the curvature tensor:
\begin{align*}
\delta a^\mu &=-\frac{1}{2}R^\mu_{\rho\lambda\alpha} a^\rho(dx^\lambda\delta x^\alpha- dx^\alpha\delta x^\lambda)\numberthis\label{Parallel transport with curvature tensor}\\
R^\mu_{\rho\lambda\alpha}&:=\Gamma^\mu_{\kappa\lambda}\Gamma^\kappa_{\rho\alpha}-\Gamma^\mu_{\kappa\alpha}\Gamma^\kappa_{\rho\lambda}+\Gamma^\mu_{\rho\lambda,\alpha}-\Gamma^\mu_{\rho\alpha,\lambda}\numberthis\label{Curvature tensor}
\end{align*}
Let $u^\mu:=x^\mu_{,\alpha}$ be a tangent vector at $P$ to a curve defined by $x^\mu=f^\mu(\alpha)$. If the transported vector is still tangent, the curve is a geodesic. Then the following equation holds (note that $\alpha$ is a simple variable in this instance):
\begin{align*}
u^\mu_{,\alpha}+\Gamma^\mu_{\rho\kappa}u^\rho u^\kappa &=0\numberthis\label{Geodesic equation}
\end{align*}
\subsection{Metric spaces and metric tensor}
A space is called metric space, if the distance between any two neighboring points for coordinates $x^\mu$ and $x^{\tilde{\mu}}$ is given by a metric tensor $g_{\mu\nu}$:
\begin{align*}
ds^2&=g_{\mu\nu}dx^\mu dx^\nu=g_{\tilde{\mu}\tilde{\nu}}dx^{\tilde{\mu}}dx^{\tilde{\nu}}\numberthis\label{Metric tensor definition}
\end{align*}
This leads to the usual transformation rule for rank 2 tensors:
\begin{align*}
g_{\mu\nu}&=x^{\tilde{\alpha}}_\mu x^{\tilde{\beta}}_\nu g_{\tilde{\alpha}\tilde{\beta}}
\end{align*}
We want the length of a vector be invariant under parallel transport, which results in the condition:
\begin{align*}
|a^\mu|^2&=g_{\mu\nu}a^\mu a^\nu\overset{!}{=}g_{\mu\nu}(P')a^{\mu}(P')a^{\nu}(P')\numberthis\label{Invariance under parallel transport}
\end{align*}
The distance between the points is $|dx^\alpha|$, so we get the approximations:
\begin{align*}
g_{\mu\nu}(P')&\simeq g_{\mu\nu}+g_{\mu\nu,\alpha}dx^\alpha\\
a^\mu(P')&=a^\mu -\Gamma^\mu_{\kappa\alpha}a^\kappa dx^\alpha
\end{align*}
Inserting these back into (\ref{Invariance under parallel transport}) and neglecting higher order terms gives:
\begin{align*}
(g_{\mu\nu,\alpha} -\Gamma^\kappa_{\nu\alpha}g_{\mu\kappa}-\Gamma^\kappa_{\mu\alpha}g_{\kappa\nu})a^\mu a^\nu dx^\alpha &=0
\end{align*}
The quantity in the brackets is just the covariant derivative $g_{\mu\nu;\alpha}$, which has to vanish for this condition to be true for any small displacements. Hence, the metric tensor is covariantly constant. This allows us to rewrite the christoffel symbol in terms of the metric tensor:
\begin{align*}
\Gamma^\mu_{\kappa\alpha}&=\frac{1}{2}g^{\mu\nu}(g_{\kappa\nu,\alpha}+g_{\nu\alpha,\kappa}-g_{\alpha\kappa,\nu})=\Gamma^\mu_{\alpha\kappa}\numberthis\label{Christoffel symbol with metric tensor}
\end{align*}
\subsection{Euler-Lagrange in Riemannian space}
Consider a curve $x^\mu(s)$ from a point $A$ to $B$ with $u^\mu :=\dot{x}^\mu (s)$. The length of such a curve is given by:
\begin{align*}
S&=\intx\int_{A}^{B}\kla\left( {g_{\mu\nu}(x^\alpha)u^\mu u^\nu} \right)^{-1/2}\,\mathrm{d}s=:\intx\int_{A}^{B}f(x^\alpha, u^\alpha)\,\mathrm{d}s
\end{align*}
Let $\epsilon\xi^\mu(s)$ be a perturbation of $x^\mu$ and $x^{\tilde{\mu}}:=x^\mu+\epsilon\xi^\mu$, where $\xi^\mu(A)=0=\xi^\mu(B)$. Note that:
\begin{align*}
f(x^{\tilde{\alpha}},u^{\tilde{\alpha}})&\simeq f(x^\alpha,u^\alpha)+\epsilon\kla\left( {\xi^\alpha f_{,\alpha}+\dot{\xi}^\alpha f_{,\dot{\alpha}}} \right)
\end{align*}
Thus, we can evaluate $\delta S:=\tilde{S}-S$:
\begin{align*}
\delta S&=\epsilon\intx\int_{A}^{B}\xi^\alpha f_{,\alpha}+(\xi^\alpha f_{,\dot{\alpha}})_{,s}-\xi^\alpha(f_{,\dot{\alpha}})_{,s}\,\mathrm{d}s=\epsilon\intx\int_{A}^{B}\klam\left[ {f_{,\alpha}-f_{,\dot{\alpha},s}} \right]\xi^\alpha \,\mathrm{d}s+\xi^\alpha f_{,\dot{\alpha}}\Big\vert_A^B\overset{!}{=}0\numberthis\label{Path length variation condition}
\end{align*}
The last term does not contribute, since the variation vanishes at the endpoints. The condition in (\ref{Path length variation condition}) then gives us the following equations:
\begin{align*}
f_{,\alpha}&=f_{,\dot{\alpha},s}\numberthis\label{Geodesic Lagrange equation}
\end{align*}
For the Lagrangian $\mathcal{L}=g_{\mu\nu}u^\mu u^\nu=f^2$ (free particle), we get $f_{,\alpha}=\frac{\mathcal{L}_{,\alpha}}{2\sqrt{\mathcal{L}}}$ and $f_{,\dot{\alpha}}=\frac{\mathcal{L}_{,\dot{\alpha}}}{2\sqrt{\mathcal{L}}}$. Thus, $\mathcal{L}_{,\alpha}=\mathcal{L}_{,\dot{\alpha},s}$. The product rule together with $u^\mu_{,\dot{\alpha}}=\delta^\mu_\alpha$ gives us $\mathcal{L}_{,\alpha}=g_{\mu\nu,\alpha}u^\mu u^\nu$ and $\mathcal{L}_{,\dot{\alpha}}=2g_{\mu\alpha}u^\mu$. Inserting these into (\ref{Geodesic Lagrange equation}) results in the geodesic equation (\ref{Geodesic equation}):
\begin{align*}
0&=\mathcal{L}_{,\alpha}-\mathcal{L}_{,\dot{\alpha},s}=g_{\mu\nu,\alpha}u^\mu u^\nu - (2g_{\mu\alpha}u^\mu)_{,s}=g_{\mu\nu,\alpha}u^\mu u^\nu - 2g_{\mu\alpha}\dot{u}^\mu - 2g_{\mu\alpha,\nu}u^\mu u^\nu
\end{align*}
Multiplying by $g^\rho_\alpha$ and using (\ref{Christoffel symbol with metric tensor}), we can rewrite this. Note that $\dot{u}^\mu =u^\mu_{,\nu}u^\nu$:
\begin{align*}
0&=g^\rho_\alpha \kla\left( {g_{\mu\alpha}\dot{u}^\mu + \frac{1}{2}\kla\left( {g_{\mu\alpha,\nu}+g_{\nu\alpha,\mu}-g_{\mu\nu,\alpha}} \right)u^\mu u^\nu} \right)=\dot{u}^\rho + \Gamma^\rho_{\mu\nu}u^\mu u^\nu=u^\rho_{;\nu}u^\nu
\end{align*}
\subsection{Riemann tensor}
The Riemann tensor can be written as:
\begin{align*}
R_{\alpha\beta\mu\nu}&=g_{\alpha\lambda}R^\lambda_{\beta\mu\nu}=\frac{1}{2}\klam\left[ {g_{\alpha\nu,\beta\mu}+g_{\beta\mu,\alpha\nu}-g_{\alpha\mu,\beta\nu}-g_{\beta\nu,\alpha\mu}+g_{\kappa\rho}\kla\left( {\Gamma^\kappa_{\alpha\nu}\Gamma^\rho_{\beta\mu}-\Gamma^\kappa_{\alpha\mu}\Gamma^\rho_{\beta\nu}} \right)} \right]\numberthis\label{Riemann tensor definition}
\end{align*}
Some properties are $R_{\alpha\beta\mu\nu}=-R_{\alpha\beta\nu\mu}$, $R_{\alpha\beta\mu\nu}=-R_{\beta\alpha\mu\nu}$ and $R_{\alpha\beta\mu\nu}=R_{\mu\nu\alpha\beta}$. The contraction of the Riemann tensor leads to the Ricci tensor:
\begin{align*}
R_{\mu\nu}&=R^\alpha_{\mu\alpha\beta}=g^{\alpha\beta}R_{\alpha\mu\beta\nu}=\Gamma^\beta_{\mu\nu,\beta}-\Gamma^\beta_{\mu\beta,\nu}+\Gamma^\beta_{\mu\nu}\Gamma^\kappa_{\kappa\beta}-\Gamma^\beta_{\mu\kappa}\Gamma^\kappa_{\nu\beta}\numberthis\label{Ricci tensor definition}
\end{align*}
Further contraction leads to the Ricci scalar:
\begin{align*}
R&=R^\mu_\mu=g^{\mu\nu}R_{\mu\nu}=g^{\mu\nu}g^{\alpha\beta}R_{\mu\nu\alpha\beta}\numberthis\label{Ricci scalar definition}
\end{align*}
With these we can define the Einstein tensor:
\begin{align*}
G_{\mu\nu}&=R_{\mu\nu}-\frac{1}{2}g_{\mu\nu}R\numberthis\label{Einstein tensor definition}
\end{align*}
It satisfies the very important condition $G^\mu_{\nu;\mu}=\kla\left( {R^\mu_\nu-\frac{1}{2}\delta^\mu_\nu R} \right)_{;\mu}=0$. Furthermore, it can be written in an expanded form:
\begin{align*}
G_{\mu\nu}&=\Gamma^\beta_{\mu\nu,\beta}-\Gamma^\beta_{\mu\beta,\nu}+\Gamma^\beta_{\mu\nu}\Gamma^\kappa_{\kappa\beta}-\Gamma^\beta_{\mu\kappa}\Gamma^\kappa_{\nu\beta}-\frac{1}{2}g_{\mu\nu}g^{\rho\lambda}\kla\left( {\Gamma^\beta_{\rho\lambda,\beta}-\Gamma^\beta_{\rho\beta,\lambda}+\Gamma^\beta_{\rho\lambda}\Gamma^\kappa_{\kappa\beta}-\Gamma^\beta_{\rho\kappa}\Gamma^\kappa_{\lambda\beta}} \right)
\end{align*}
\subsection{Examples}
\subsubsection{Geodesics on a sphere}
Consider the surface of a sphere with the standard metric $ds^2=d\theta^2+\sin^2\theta d\phi^2$ and Christoffel symbols $\Gamma^1_{22}=-\sin\theta\cos\theta$ and $\Gamma^2_{12}=\cot\theta$. The geodesic equations for $\theta(t)$ and $\phi(t)$ with $u^k=x^k_{,t}(t)$ are given by:
\begin{align*}
0&=u^k_{;j}u^j=x^k_{,tt}+\Gamma^k_{ij}x^i_{,t}x^j_{,t}\\
\ddot{\phi}&=-2\cot\theta\dot{\theta}\dot{\phi}\ \ \Rightarrow\ \ \log\dot{\phi}=\log a-2\log\sin\theta\ \ \Rightarrow\ \ \dot{\phi}=a\csc^2\theta\\
\ddot{\theta}&=\sin\theta\cos\theta\dot{\phi}^2=a^2\csc^2\theta\cot\theta\ \ \Rightarrow\ \ \dot{\theta}^2=b^2-a^2\csc^2\theta
\end{align*}
Choose $b=a$ and use seperation of variables to solve the last equation:
\begin{align*}
iat+d&=\intx\int_{ }^{ }\tan\theta\,\mathrm{d}\theta=-\log\cos\theta\\
\cos\theta(t)&=e^{iat}\cos\theta_0\equiv\cos at\cos\theta_0
\end{align*}
Let $c=\cos\theta_0$ and solve for $\phi(t)$. We can always choose $\phi_0=0$:
\begin{align*}
\phi(t)+b&=a\intx\int_{ }^{ }\frac{1}{1-c^2\cos^2at}\,\mathrm{d}t=\frac{1}{\sqrt{1-c^2}}\arctan\frac{\tan at}{\sqrt{1-c^2}}
\end{align*}
\section{Variational approach}
\subsection{Lagrange and Hamiltonian}
Let us compare the structure of Lagrange mechanics in the discrete and continuous case:
\begin{table}[h!]
\centering
\begin{tabular}{c|c|c|c}
&Particles&Fields&4-Fields\\ \hline
System state&$q_j(t)$&$\psi(x^\mu)$&$A^\nu(x^\mu)$\\ \hline
Lagrangian&$L(q_j,\dot{q}_j)$&$\mathcal{L}(\psi,\psi_{,\mu})$&$\mathcal{L}(A^\nu,A^\nu_{,\mu})$\\ \hline
Equation&$L_{,q}=L_{,\dot{q},t}$&$\mathcal{L}_{,\psi}=\mathcal{L}_{,\psi_{,\mu},\mu}$&$\mathcal{L}_{,A^\nu}=\mathcal{L}_{,A^\nu_{,\mu},\mu}$\\ \hline
Hamiltonian&$\sum_j\dot{q}_jL_{,\dot{q_j}}-L$&$\intr\int_{\mathbb{R}^{3}}{\psi_{,0}\mathcal{L}_{,\psi_{,0}}-\mathcal{L}}\,\mathrm{d}^{3}x$&
\end{tabular}
\caption{Lagrangian for particles and fields}
\label{Lagrangian for particles and fields}
\end{table}
\subsection{Energy momentum tensor}
The integrand of the Hamiltonian can be considered as an energy density $t^0_0=\psi_{,0}\mathcal{L}_{,\psi_{,0}}-\mathcal{L}$. This leads to an educated guess for the structure of the full energy momentum tensor:
\begin{align*}
t^\mu_\nu&=\psi_{,\nu}\mathcal{L}_{,\psi_{,\mu}}-\delta^\mu_\nu\mathcal{L}\numberthis\label{Energy momentum tensor definition}
\end{align*}
This tensor should fulfill the condition $t^\mu_{\nu,\mu}=0$. Note that $\mathcal{L}_{,\mu}=\psi_{,\mu}\mathcal{L}_{,\psi}+\psi_{,\alpha\mu}\mathcal{L}_{,\psi_{,\alpha}}$:
\begin{align*}
t^\mu_{\nu,\mu}&=\kla\left( {\psi_{,\nu}\mathcal{L}_{,\psi_{,\mu}}} \right)_{,\mu}-\delta^\mu_\nu\mathcal{L}_{,\mu}=\psi_{,\nu}\mathcal{L}_{,\psi_{,\mu},\mu}-\psi_{,\nu}\mathcal{L}_{,\psi}\overset{(\ref{Lagrangian for particles and fields})}{=}0
\end{align*}
\subsection{Scalar fields}
Consider the one-component field $\psi(x^\mu)$ with the Lagrange density $\mathcal{L}=\frac{1}{2}\kla\left( {\eta^{\mu\nu}\psi_{,\mu}\psi_{,\nu}-m^2\psi^2} \right)$ for a constant $m$. The field equation then becomes:
\begin{align*}
0&=\mathcal{L}_{,\psi_{,\mu},\mu}-\mathcal{L}_{,\psi}=(\eta^{\mu\nu}\psi_{,\nu})_{,\mu}+m^2\psi=\psi^{,\mu}_{,\mu}+m^2\psi=(\square +m^2)\psi\numberthis\label{Klein Gordon equation}
\end{align*}
From (\ref{Energy momentum tensor definition}) we can derive the energy momentum tensor for this case:
\begin{align*}
t^\mu_\nu&=\psi_{,\nu}\eta^{\mu\alpha}\psi_{,\alpha}-\delta^\mu_\nu\mathcal{L}=\psi_{,\nu}\psi^{,\mu}-\frac{1}{2}\delta^\mu_\nu\kla\left( {\psi_{,\alpha}\psi^{,\alpha}-m^2\psi^2} \right)
\end{align*}
The energy density is then given by $2t^0_0=(\partial_t\psi)^2-(\nabla\psi)^2+m^2\psi^2$
\subsection{Multi component fields}
Define the EM field tensor $F^{\mu\nu}=A^{\mu,\nu}-A^{\nu,\mu}$. Then the Lagrange formalism gives us the Maxwell equations in free space if we vary the Lagrange density $\mathcal{L}:=-\frac{1}{16\pi}F_{\mu\nu}F^{\mu\nu}$ with respect to the 4-vector $A^\nu$.\\
A similar approach yields the equations for the gravitational field in linear approximation, if we vary $\mathcal{L}:=\frac{1}{4}\kla\left( h_{\mu\nu,\lambda}h^{\mu\nu,\lambda}-2h^{,\mu}_{\mu\nu}h^{\nu\lambda}_{,\lambda}+2h^{,\mu}_{\mu\nu}h^{,\nu}-h_{,\mu}h^{,\mu} \right)$ with respect to $h^{\alpha\beta}$.
\subsection{Nonlinear Einstein equations for gravity}
Define the Lagrange density $\mathcal{L}:=\frac{1}{16\pi}R\sqrt{-g}$ with $G,c=1$. The curvature scalar $R$ can be written as follows:
\begin{align*}
R\sqrt{-g}&=\kla\left( {g^{\mu\nu}\sqrt{-g}[\delta^\alpha_\mu\Gamma^\beta_{\beta\nu}-\Gamma^\alpha_{\mu\nu} ]} \right)_{,\alpha}+g^{\mu\nu}\sqrt{-g}\kla\left( {\Gamma^\alpha_{\mu\beta}\Gamma^\beta_{\alpha\nu}-\Gamma^\alpha_{\mu\nu}\Gamma^\beta_{\beta\alpha}} \right)\numberthis\label{Curvature Scalar}
\end{align*}
Since the first term is a divergence, its contribution to the action via $S=\int\mathcal{L}$ can be converted to an integral over the boundary of the volume. But in the variational approach, $\delta g_{\mu\nu}$ vanishes at the boundary, which means that it does not effect the variation $\delta S$ in any way. Hence, the equations of motion can be derived from an effective Lagrangian density only:
\begin{align*}
\mathcal{L}_e&\propto g^{\mu\nu}\sqrt{-g}\kla\left( {\Gamma^\alpha_{\mu\beta}\Gamma^\beta_{\alpha\nu}-\Gamma^\alpha_{\mu\nu}\Gamma^\beta_{\beta\alpha}} \right)\numberthis\label{Effective Lagrange density nonlinear gravity}
\end{align*}
For the derivation of the equations we will need a few tools. Let $D$ be a diagonal matrix and define $A:=e^D$, which is also diagonal. Then we get $\det A=e^{\tecxt\text{tr\,}}D$ and thus follows:
\begin{align*}
\log\det A&=\tecxt\text{tr\,}\log A\numberthis\label{log det A equals trace log A}
\end{align*}
Then we can calculate the derivative of the determinant by taking an arbitrary differential $\partial$ on both sides of (\ref{log det A equals trace log A}):
\begin{align*}
\frac{\partial (\det A)}{\det A}&=\tecxt\text{tr\,}\frac{\partial A}{A}=\tecxt\text{tr\,}A^{-1}\partial A
\end{align*}
Now letting $A=g_{\mu\nu}$ and $g=\det g_{\mu\nu}$ with $A^{-1}=g^{\mu\nu}$ yields:
\begin{align*}
g_{,\alpha}&=g\,\tecxt\text{tr\,}g^{\mu\nu}g_{\mu\nu,\alpha}=gg^{\mu\nu}g_{\mu\nu,\alpha}\\
\sqrt{-g}_{,g_{\mu\nu}}&=\frac{1}{2}\sqrt{-g} g^{\mu\nu}
\end{align*} 
Recall the Christoffel symbol $\Gamma^\mu_{\alpha\beta}=\frac{1}{2}g^{\mu\kappa}(g_{\alpha\kappa,\beta}+g_{\beta\kappa,\alpha}-g_{\alpha\beta,\kappa})$ and write out (\ref{Effective Lagrange density nonlinear gravity}):
\begin{align*}
\Gamma^\mu_{\alpha\beta,g_{\lambda\rho,\nu}}&=\frac{1}{2}g^{\mu\kappa}(\delta^{\lambda\rho\nu}_{\alpha\kappa\beta}+\delta^{\lambda\rho\nu}_{\beta\kappa\alpha}-\delta^{\lambda\rho\nu}_{\alpha\beta\kappa})=\frac{1}{2}g^{\mu\rho}(\delta^{\lambda\nu}_{\alpha\beta}+\delta^{\lambda\nu}_{\beta\alpha})-\frac{1}{2}g^{\mu\nu}\delta^{\lambda\rho}_{\alpha\beta}
\end{align*}
test
\begin{align*}
\mathcal{L}&=\sqrt{-g}g^{\mu\nu}g^{\alpha\kappa}g^{\beta\lambda}\big[ (g_{\mu\kappa,\beta}+g_{\beta\kappa,\mu}-g_{\beta\mu,\kappa})(g_{\alpha\lambda,\nu}+g_{\nu\lambda,\alpha}-g_{\alpha\nu,\lambda}) -\\
&\hspace{3cm}-(g_{\mu\kappa,\nu}+g_{\nu\kappa,\mu}-g_{\mu\nu,\kappa})(g_{\alpha\lambda,\beta}+g_{\beta\lambda,\alpha}-g_{\alpha\beta,\lambda})\big]\\
&=\sqrt{-g}g^{\mu\nu}g^{\alpha\kappa}g^{\beta\lambda}\big[g_{\mu\kappa,\beta}g_{\alpha\lambda,\nu}+g_{\beta\kappa,\mu}g_{\alpha\lambda,\nu}-g_{\beta\mu,\kappa}g_{\alpha\lambda,\nu}+g_{\mu\kappa,\beta}g_{\nu\lambda,\alpha}+g_{\beta\kappa,\mu}g_{\nu\lambda,\alpha}-g_{\beta\mu,\kappa}g_{\nu\lambda,\alpha}\\
&\hspace{3cm} -g_{\mu\kappa,\beta}g_{\alpha\nu,\lambda}-g_{\beta\kappa,\mu}g_{\alpha\nu,\lambda}+g_{\beta\mu,\kappa}g_{\alpha\nu,\lambda}-g_{\mu\kappa,\nu}g_{\alpha\lambda,\beta}-g_{\nu\kappa,\mu}g_{\alpha\lambda,\beta}+g_{\mu\nu,\kappa}g_{\alpha\lambda,\beta}\\
&\hspace{3cm} -g_{\mu\kappa,\nu}g_{\beta\lambda,\alpha}-g_{\nu\kappa,\mu}g_{\beta\lambda,\alpha}+g_{\mu\nu,\kappa}g_{\beta\lambda,\alpha}+g_{\mu\kappa,\nu}g_{\alpha\beta,\lambda}+g_{\nu\kappa,\mu}g_{\alpha\beta,\lambda}-g_{\mu\nu,\kappa}g_{\alpha\beta,\lambda}\big]
\end{align*}
test
\begin{align*}
0&=\mathcal{L}_{,g_{\mu\nu,\alpha},\alpha}-\mathcal{L}_{,g_{\mu\nu}}\\
\frac{\mathcal{L}_{,g_{ij,k}}}{\sqrt{-g}}&=g^{i\nu}g^{j\alpha}g^{k\lambda}g_{\alpha\lambda,\nu}+g^{k}g^{i}g^{j}g_{\alpha\lambda,\nu}-g^{j}g^{k}g^{i}g_{\alpha\lambda,\nu}+g^{}g^{}g^{}g_{\nu\lambda,\alpha}+g^{}g^{}g^{}g_{\nu\lambda,\alpha}-g^{}g^{}g^{}g_{\nu\lambda,\alpha}
\end{align*}
test2
\begin{align*}
\mathcal{L}_{,g_{\kappa\rho}}&=-\frac{1}{2} g^{\kappa\rho}\mathcal{L}\\
\mathcal{L}_{,g_{\kappa\rho,\lambda}}&=\frac{1}{2}\sqrt{-g}g^{\mu\nu}\kla\left( {\Gamma^\beta_{\alpha\nu}[g^{\alpha\rho}(\delta^{\kappa\lambda}_{\mu\beta}+\delta^{\kappa\lambda}_{\beta\mu})-g^{\alpha\lambda}\delta^{\kappa\rho}_{\mu\beta}]+\Gamma^\alpha_{\mu\beta}[g^{\beta\rho}(\delta^{\kappa\lambda}_{\nu\alpha}+\delta^{\kappa\lambda}_{\alpha\nu})-g^{\beta\lambda}\delta^{\kappa\rho}_{\nu\alpha}]-\dots} \right)=\\
&=\frac{1}{2}\sqrt{-g}\kla\left( {g^{\alpha\rho}[g^{\kappa\nu}\Gamma^\lambda_{\alpha\nu}+g^{\lambda\nu}\Gamma^\kappa_{\alpha\nu}]-g^{\alpha\lambda}g^{\kappa\nu}\Gamma^\rho_{\alpha\nu}+g^{\beta\rho}[g^{\mu\kappa}\Gamma^\lambda_{\mu\beta}+g^{\mu\lambda}\Gamma^\kappa_{\mu\beta}]-g^{\mu\kappa}g^{\beta\lambda}\Gamma^\rho_{\mu\beta}} \right)
\end{align*}
%\Gamma^\alpha_{\mu\beta}\Gamma^\beta_{\alpha\nu}-\Gamma^\alpha_{\mu\nu}\Gamma^\beta_{\beta\alpha}







































\newpage
1

\end{document}